% from ega4-21c-fr.tex
\oldpage[IV-4]{312}
$f_0:X_0\to Y_0$ tels que $g$, $h$ et $f$ se déduisent de $g_0$, $h_0$, $f_0$ par le
changement de base $S\to S_0$ ((8.9.1) et (5.13.3)); on peut en outre supposer $g_0$
plat (11.2.7), $h_0$ lisse (17.7.9) et $f_0$ propre (8.10.5, (xii)). D'autre part, la
projection sur $S_0$ du point générique $\eta$ de $S$ est le point générique $\eta_0$ de
$S_0$ et il résulte de (2.7.1, (viii)) que le morphisme
$(f_0)_{\eta_0}:(X_0)_{\eta_0}\to(Y_0)_{\eta_0}$ est un isomorphisme; enfin le point
fermé $s$ de $S$ est au-dessus du point fermé $s_0$ de $S_0$, donc la fibre
$(X_0)_{s_0}$ de $g_0$ est géométriquement irréductible (4.5.6).

On voit donc qu'on peut se borner au cas où $S$ est le spectre d'un anneau local intègre
noethérien, de point fermé $s$, affaiblir l'hypothèse sur les fibres en supposant seulement
que $X_s$ et $X_\eta$ sont géométriquement irréductibles, et prouver que sous ces
hypothèses $f_s$ est une immersion ouverte. Il y a alors un anneau de valuation discrète
$A'$ et un morphisme $S'=\operatorname{Spec}(A')\to S$ qui transforme le point fermé
$a$ de $S'$ en $s$ et le point générique $b$ de $S'$ en $\eta$ $(\mathbf{II}, 7.1.9)$;
soient $g':X'\to S'$, $h':Y'\to S'$, $f':X'\to Y'$ les morphismes déduits de $g$,
$h$, $f$ par le changement de base $S'\to S$, qui vérifient de nouveau les hypothèses
vérifiées par $g$, $h$, et $f$ dans l'énoncé (21.12.15); si l'on prouve que
$f'_a:X'_a\to Y'_a$ est une immersion ouverte, il résultera de (2.7.1, (x)) qu'il en
sera de même de $f_s:X_s\to Y_s$. On peut donc finalement se borner au cas où $S$ est
le \emph{spectre d'un anneau de valuation discrète}.

Alors, puisque $h:Y\to S$ est lisse, $Y$ est régulier (17.5.8), et puisque la question
est locale sur $Y$, on peut se borner au cas où $Y$ est intègre. Comme $g:X\to S$ est
plat, les points maximaux de $X$ sont contenus dans $X_\eta$ (2.3.4) et puisque
$f_\eta$ est un isomorphisme, $X$ est irréductible et l'anneau local en son point
générique est un corps; de plus, par platitude (3.3.2), on voit que $X$ n'a pas de cycles
premiers immergés, donc $X$ est réduit (3.2.1) et par suite intègre et $f$ est un
morphisme birationnel et séparé. Pour prouver que $f$ est une immersion ouverte, on peut
donc appliquer le critère (21.12.13); pour voir que $f$ est quasi-fini aux points de
$X^{(1)}$, on remarque que l'assertion est évidente en ceux de ces points qui
appartiennent à $X_\eta$; le seul point de $X^{(1)}$ appartenant à $X_s$ est le point
générique $x$ de $X_s$, en vertu de (6.1.1). Or, il résulte de (2.4.6) et de (14.2.2)
que les morphismes $g$ et $h$ sont équidimensionnels; comme $X_\eta$ et $Y_\eta$ sont
isomorphes, les composantes irréductibles de $X_s$ et de $Y_s$ ont toutes même
dimension. Mais par hypothèse $X_s$ est irréductible et on a vu que $f$ est surjectif,
donc il en est de même de $f_s:X_s\to Y_s$, ce qui entraîne que $Y_s$ est aussi
irréductible; si $y$ est le point générique de $Y_s$, on a donc $f(x)=y$, et la relation
$\dim(X_s)=\dim(Y_s)$ entraîne alors, en vertu de (5.6.6), que
$\dim(f_s^{-1}(y))=0$, autrement dit $f_s$ (et par suite aussi $f$) est bien quasi-fini
au point $x$, ce qui achève la démonstration.
\end{proof}

\begin{corollary}[21.12.16]
\label{IV.21.12.16}
Soient $g:X\to S$ un morphisme propre, plat et de présentation finie, $h:Y\to S$ un
morphisme lisse, propre et de présentation finie, $f:X\to Y$ un $S$-morphisme. On
suppose que les fibres $X_s=g^{-1}(s)$ de $g$ sont géométriquement irréductibles. Soit
$U$ l'ensemble des $s\in S$ tels que $f_s:X_s\to Y_s$ soit un isomorphisme. Alors
$U$ est ouvert et fermé dans $S$, et le morphisme $g^{-1}(U)\to h^{-1}(U)$,
restriction de $f$, est un isomorphisme.
\end{corollary}

\begin{proof}
La dernière assertion résultera de la première et de (17.9.5). On sait déjà (9.6.1,
(xi)) que $U$ est localement constructible dans $S$. En vertu de (1.10.1), il suffit de
prouver que $U$ est stable par \emph{spécialisation} et par \emph{générisation}. Pour
démontrer ces assertions on
